\chapter{Desarrollo}
\label{chap:desarrollo}

\lettrine{E}{l} desarrollo del sistema se ha realizado siguiendo una metodología incremental-iterativa que ha permitido evolucionar desde prototipos funcionales simples hasta una arquitectura modular robusta. Este capítulo documenta el proceso de desarrollo completo aplicando de forma sistemática las fases de análisis, diseño e implementación en cada iteración, describiendo los artefactos generados, las decisiones arquitectónicas adoptadas y las validaciones realizadas.

Cada iteración sigue el ciclo metodológico establecido: análisis de requisitos específicos, diseño de la solución, implementación de componentes, y validación mediante pruebas funcionales. Esta aproximación sistemática ha resultado fundamental para gestionar la complejidad inherente del desarrollo de sistemas \acrshort{bci} integrados.

\section{Análisis de requisitos del sistema}
\label{sec:analisis_requisitos_sistema}

El análisis de requisitos se realizó de forma evolutiva durante las primeras iteraciones, refinándose progresivamente a medida que se profundizaba en el dominio del problema. La identificación y clasificación sistemática de requisitos proporcionó la base para el diseño arquitectónico y la planificación de iteraciones.

\subsection{Requisitos funcionales}
\label{subsec:requisitos_funcionales}

Los requisitos funcionales del sistema se clasificaron y documentaron de forma sistemática para guiar el desarrollo incremental:

\begin{itemize}
    \item \textbf{RF-01}: Captura de eventos de teclado con precisión temporal de milisegundos
    \item \textbf{RF-02}: Captura de eventos de ratón con múltiples modos de \textit{tracking}
    \item \textbf{RF-03}: Soporte para controladores \acrshort{ps4} mediante \textit{SDL2}
    \item \textbf{RF-04}: Adquisición de señales \acrshort{eeg} mediante BrainFlow
    \item \textbf{RF-05}: Soporte para múltiples dispositivos \acrshort{eeg} comerciales
    \item \textbf{RF-06}: Procesamiento de señales en tiempo real con algoritmos avanzados
    \item \textbf{RF-07}: Visualización en tiempo real de \gls{biosenial}es y eventos
    \item \textbf{RF-08}: Sincronización temporal precisa entre fuentes de datos
    \item \textbf{RF-09}: Exportación a formatos estándar (\acrshort{bids}, \acrshort{edf}+, \acrshort{hdf5}, \acrshort{csv})
    \item \textbf{RF-10}: Interfaces múltiples (\acrshort{gui}, \acrshort{cli}, \acrshort{api} web)
    \item \textbf{RF-11}: Sistema de ayuda contextual integrado
    \item \textbf{RF-12}: Monitoreo de calidad de señal en tiempo real
\end{itemize}

\subsection{Requisitos no funcionales}
\label{subsec:requisitos_no_funcionales}

Los requisitos no funcionales establecieron las restricciones de calidad y rendimiento del sistema:

\begin{itemize}
    \item \textbf{RNF-01}: Latencia de captura de eventos inferior a 50ms
    \item \textbf{RNF-02}: Soporte para frecuencias de muestreo \acrshort{eeg} hasta 1000Hz
    \item \textbf{RNF-03}: Arquitectura extensible basada en principios \acrshort{solid}
    \item \textbf{RNF-04}: Operaciones \textit{thread-safe} para procesamiento concurrente
    \item \textbf{RNF-05}: Interfaces intuitivas para usuarios no técnicos
    \item \textbf{RNF-06}: Compatibilidad multiplataforma (\textit{Windows}, \textit{Linux}, \textit{macOS})
    \item \textbf{RNF-07}: Tipado estático completo con \textit{type hints}
    \item \textbf{RNF-08}: Manejo robusto de errores con recuperación automática
    \item \textbf{RNF-09}: Documentación técnica completa y actualizada
    \item \textbf{RNF-10}: Cumplimiento de estándares de codificación establecidos
\end{itemize}

\subsection{Casos de uso principales}
\label{subsec:casos_uso_principales}

Los casos de uso identificados guiaron el desarrollo de funcionalidades específicas:

\begin{itemize}
    \item \textbf{CU-01}: Generación de \textit{datasets} para entrenamiento \acrshort{ml}
    \item \textbf{CU-02}: Análisis de señales en tiempo real durante experimentos
    \item \textbf{CU-03}: Desarrollo de aplicaciones \acrshort{bci} personalizadas
    \item \textbf{CU-04}: Investigación académica en neurociencia
    \item \textbf{CU-05}: Validación de paradigmas experimentales
    \item \textbf{CU-06}: Integración con herramientas de análisis existentes
\end{itemize}

\section{Primera iteración: Módulo de captura de entrada}
\label{sec:iteracion_1}

La primera iteración se enfocó en desarrollar las capacidades básicas de captura de eventos de entrada del usuario, estableciendo los fundamentos técnicos para el procesamiento temporal preciso.

\subsection{Análisis de la iteración}
\label{subsec:analisis_iteracion_1}

\textbf{Objetivos específicos}: Implementar captura básica de eventos de teclado y ratón con precisión temporal (RF-01, RF-02).

\textbf{Requisitos abordados}: RF-01, RF-02, RNF-01, RNF-06

\textbf{Casos de uso}: CU-01 (parcial), CU-02 (parcial)

\textbf{Criterios de aceptación}:
\begin{itemize}
    \item Captura de eventos de teclado con \textit{timestamps} precisos
    \item Filtrado configurable de teclas específicas
    \item Múltiples modos de \textit{tracking} para ratón
    \item Exportación básica en formato \acrshort{json}
\end{itemize}

\subsection{Diseño de la iteración}
\label{subsec:diseno_iteracion_1}

\textbf{Arquitectura inicial}: Diseño monolítico con captura directa mediante \textit{pynput}

\textbf{Componentes principales}:
\begin{itemize}
    \item Capturador de eventos de teclado
    \item Capturador de eventos de ratón con modos múltiples
    \item Sistema de \textit{timestamping} de alta precisión
    \item Exportador \acrshort{json} básico
\end{itemize}

\textbf{Decisiones de diseño}:
\begin{itemize}
    \item Uso de \textit{pynput} para compatibilidad multiplataforma
    \item Implementación de múltiples modos de \textit{tracking} de ratón: \textit{delta}, ángulo, eje, \textit{joystick}
    \item Sistema de filtrado configurable para teclas específicas
\end{itemize}

\subsection{Implementación de la iteración}
\label{subsec:implementacion_iteracion_1}

\textbf{Tecnologías utilizadas}: \textit{Python 3.8+}, \textit{pynput}, \textit{time} (para \textit{timestamps})

\textbf{Desarrollo del \textit{script} básico}:
Se desarrolló un \textit{script} inicial utilizando \textit{pynput} para registrar eventos de teclado con \textit{timestamps} precisos. Una vez verificado el funcionamiento básico, se añadió un sistema de filtrado que permitía capturar únicamente teclas preestablecidas específicas para evitar ruido en los datos.

\textbf{Soporte para eventos de ratón}:
La incorporación de eventos de ratón presentó complejidad específica en aplicaciones donde el cursor se mantiene bloqueado en el centro de la pantalla. Se implementaron cuatro modos de representación del movimiento:
\begin{itemize}
    \item \textit{Delta}: Movimiento en bruto con direcciones y magnitudes
    \item \textit{Joystick}: Representación analógica del movimiento
    \item Ángulo: Direcciones en 8 puntos cardinales con velocidad por eje
    \item Eje: Movimiento separado por componentes X e Y
\end{itemize}

\textbf{Gestión de eventos de liberación}:
Se implementó filtrado configurable para eventos de liberación (\textit{release}) de teclas y botones de ratón, permitiendo capturar solo los eventos relevantes para cada experimento específico.

\textbf{Migración tecnológica fallida}:
Se intentó migrar el \textit{script} a \textit{Julia} por su rendimiento superior y tipado fuerte, pero la precompilación requerida en cada ejecución y la dificultad para depurar errores complejos motivaron continuar con \textit{Python} debido a la naturaleza experimental iterativa del desarrollo.

\textbf{Implementación de \textit{type hints}}:
Se añadió tipado estático básico al \textit{script} para mejorar la robustez y facilitar el mantenimiento del código.

\subsection{Validación de la iteración}
\label{subsec:validacion_iteracion_1}

\textbf{Pruebas realizadas}:
\begin{itemize}
    \item Verificación de precisión temporal en \textit{timestamps}
    \item Validación de filtrado de teclas específicas
    \item Pruebas de todos los modos de \textit{tracking} de ratón
    \item Verificación de exportación \acrshort{json} correcta
\end{itemize}

\textbf{Resultados}: Cumplimiento satisfactorio de RF-01, RF-02 y RNF-01. Latencias medidas < 10ms para eventos de entrada.

\section{Segunda iteración: Módulo de procesamiento \acrshort{eeg}}
\label{sec:iteracion_2}

La segunda iteración desarrolló las capacidades de captura y procesamiento de señales \acrshort{eeg}, estableciendo la base para el análisis de \gls{biosenial}es en tiempo real.

\subsection{Análisis de la iteración}
\label{subsec:analisis_iteracion_2}

\textbf{Objetivos específicos}: Implementar adquisición de señales \acrshort{eeg} con visualización en tiempo real (RF-04, RF-05, RF-07).

\textbf{Requisitos abordados}: RF-04, RF-05, RF-06, RF-07, RF-11, RF-12, RNF-02, RNF-05

\textbf{Casos de uso}: CU-02, CU-04, CU-05

\textbf{Criterios de aceptación}:
\begin{itemize}
    \item Conectividad con dispositivos BrainFlow
    \item Visualización en tiempo real de múltiples canales
    \item Monitoreo de calidad de señal
    \item Sistema de ayuda contextual
    \item Grabación en formato \acrshort{csv}
\end{itemize}

\subsection{Diseño de la iteración}
\label{subsec:diseno_iteracion_2}

\textbf{Arquitectura}: Integración con BrainFlow mediante patrón \textit{Adapter}

\textbf{Componentes principales}:
\begin{itemize}
    \item Adaptador BrainFlow para abstracción de dispositivos
    \item Visualizador multicanal con detección automática
    \item Monitor de calidad de señal (\textit{railed percentage})
    \item Visualizador topográfico (\textit{Brain Map})
    \item Sistema de ayuda contextual
\end{itemize}

\textbf{Decisiones de diseño}:
\begin{itemize}
    \item Migración a \textit{PyQt5} por rendimiento superior a \textit{tkinter}
    \item Implementación de sistema 10-20 para posicionamiento de electrodos
    \item Integración de métricas de estado mental en tiempo real
\end{itemize}

\subsection{Implementación de la iteración}
\label{subsec:implementacion_iteracion_2}

\textbf{Tecnologías utilizadas}: \textit{PyQt5}, BrainFlow, \textit{NumPy}, \textit{matplotlib}

\textbf{Experimentación inicial con BrainFlow}:
Se comenzó utilizando la placa sintética de BrainFlow para validar la \acrshort{api} y establecer patrones de adquisición de datos. Se implementó un grabador básico para registrar datos en archivos, validando el flujo completo de adquisición, procesamiento y almacenamiento.

\textbf{Migración de interfaz gráfica}:
La implementación inicial en \textit{tkinter} mostró limitaciones de rendimiento para visualización en tiempo real y falta de \textit{widgets} especializados. La migración a \textit{PyQt5} proporcionó mejor rendimiento y el sistema de señales/\textit{slots} para comunicación entre componentes.

\textbf{Desarrollo del sistema de ayuda}:
Se implementaron más de 25 \textit{widgets} de ayuda contextual que cubren desde conceptos básicos de neurofisiología hasta parámetros técnicos específicos. La información se compiló de documentación oficial de BrainFlow y recursos educativos \textit{online}.

\textbf{Visualizador topográfico (\textit{Brain Map})}:
Se implementó un visualizador que representa la distribución de electrodos según el sistema 10-20 dividido en 4 regiones anatómicas (frontal, central, parietal, occipital). Integra indicadores de calidad de señal por electrodo utilizando codificación de colores.

\textbf{Grabador \acrshort{csv} con metadatos}:
Se desarrolló un sistema de grabación que incluye metadatos en las cabeceras: fecha/hora, identificación de placa, \textit{preset}, número de canales, duración, y columnas para \textit{timestamps}, marcadores y canal de control BrainFlow.

\subsection{Validación de la iteración}
\label{subsec:validacion_iteracion_2}

\textbf{Pruebas realizadas}:
\begin{itemize}
    \item Conectividad con placa sintética BrainFlow
    \item Validación de visualización multicanal
    \item Verificación de calidad de señal
    \item Pruebas de grabación con metadatos
\end{itemize}

\textbf{Resultados}: Cumplimiento de RF-04, RF-05, RF-07, RF-11, RF-12. Rendimiento validado para frecuencias hasta 250Hz.

\section{Tercera iteración: Primer intento de integración}
\label{sec:iteracion_3}

Esta iteración abordó la integración visual y funcional de ambos módulos, revelando limitaciones arquitectónicas que motivarían la refactorización posterior.

\subsection{Análisis de la iteración}
\label{subsec:analisis_iteracion_3}

\textbf{Objetivos específicos}: Integrar módulos existentes con sincronización temporal precisa (RF-08).

\textbf{Requisitos abordados}: RF-08, RF-09, RNF-04

\textbf{Casos de uso}: CU-01, CU-02

\textbf{Limitaciones identificadas}:
\begin{itemize}
    \item Arquitectura monolítica dificultaba integración
    \item Falta de sincronización temporal entre módulos
    \item Acoplamiento alto impedía modularidad
\end{itemize}

\subsection{Diseño de la iteración}
\label{subsec:diseno_iteracion_3}

\textbf{Aproximación inicial}: Integración visual mediante coordinador central

\textbf{Tecnología de sincronización}: \textit{Lab Streaming Layer} (\acrshort{lsl})

\textbf{Problemas de diseño identificados}:
\begin{itemize}
    \item Archivos monolíticos de gran tamaño
    \item Acoplamiento estrecho entre componentes
    \item Dificultad para identificar puntos de integración
\end{itemize}

\subsection{Implementación de la iteración}
\label{subsec:implementacion_iteracion_3}

\textbf{Integración visual exitosa}:
Se logró integrar ambos módulos en una interfaz unificada que permitía controlar la captura de señales \acrshort{eeg} y eventos de entrada simultáneamente. Se implementó coordinación de grabadores para inicio/parada sincronizada.

\textbf{Unificación de formatos}:
Se añadió soporte \acrshort{csv} al módulo de captura de entrada (que previamente solo soportaba \acrshort{json}) para facilitar la combinación de datos de ambos módulos.

\textbf{Problemas con \acrshort{lsl}}:
La integración \acrshort{lsl} resultó inviable debido a la estructura monolítica de los módulos y archivos de gran tamaño. La complejidad de depuración en código fuertemente acoplado dificultó identificar y corregir problemas de integración.

\textbf{Problemas de sincronización}:
Cada módulo utilizaba su propio sistema de \textit{timestamping} sin coordinación externa, impidiendo la sincronización temporal precisa requerida para aplicaciones \acrshort{bci}.

\subsection{Validación de la iteración}
\label{subsec:validacion_iteracion_3}

\textbf{Resultado}: Iteración fallida que motivó refactorización arquitectónica fundamental.

\textbf{Lecciones aprendidas}:
\begin{itemize}
    \item Necesidad de arquitectura modular
    \item Importancia de principios \acrshort{solid}
    \item Requisito de interfaces bien definidas
\end{itemize}

\section{Cuarta iteración: Refactorización \acrshort{solid}}
\label{sec:iteracion_4}

La refactorización completa de ambos módulos siguiendo principios \acrshort{solid} estableció una base arquitectónica para futuras integraciones.

\subsection{Análisis de la iteración}
\label{subsec:analisis_iteracion_4}

\textbf{Objetivos específicos}: Reestructurar código existente aplicando principios \acrshort{solid} (RNF-03).

\textbf{Requisitos abordados}: RNF-03, RNF-04, RNF-08, RNF-09

\textbf{Principios aplicados}:
\begin{itemize}
    \item \textit{Single Responsibility}: Cada clase con una responsabilidad específica
    \item \textit{Open-Closed}: Extensibilidad mediante interfaces y abstracciones
    \item \textit{Liskov Substitution}: Implementaciones intercambiables
    \item \textit{Interface Segregation}: Interfaces específicas para cada funcionalidad
    \item \textit{Dependency Inversion}: Dependencias mediante abstracciones
\end{itemize}

\subsection{Diseño de la iteración}
\label{subsec:diseno_iteracion_4}

\textbf{Nueva arquitectura}: Diseño basado en servicios \textit{singleton} y patrones dirigidos por eventos

\textbf{Componentes arquitectónicos}:
\begin{itemize}
    \item Servicios centralizados (\textit{CollectorService}, \textit{DataAccessService})
    \item Interfaces abstractas para extensibilidad
    \item Gestores especializados por funcionalidad
    \item Sistema de \textit{callbacks} interno
\end{itemize}

\textbf{Patrones de diseño implementados}:
\begin{itemize}
    \item \textit{Singleton} para servicios centrales
    \item \textit{Factory} para creación de colectores
    \item \textit{Observer} para comunicación entre componentes
    \item \textit{Strategy} para diferentes algoritmos de procesamiento
\end{itemize}

\subsection{Implementación de la iteración}
\label{subsec:implementacion_iteracion_4}

\textbf{Refactorización del módulo de entrada}:
Se separaron responsabilidades específicas en clases independientes: gestores de colectores, servicios de acceso a datos, adaptadores para comunicación. La granularidad mejoró la legibilidad y mantenibilidad del código.

\textbf{Refactorización del módulo \acrshort{eeg}}:
Se implementaron servicios \textit{singleton} \textit{thread-safe} para recursos compartidos, separación procesamiento-presentación, y sistema de \textit{callbacks} internos para comunicación desacoplada.

\textbf{Validación funcional}: Se realizaron pruebas manuales exhaustivas después de cada cambio para garantizar preservación completa de funcionalidad existente.

\subsection{Validación de la iteración}
\label{subsec:validacion_iteracion_4}

\textbf{Mejoras observadas}:
\begin{itemize}
    \item Archivos de menor tamaño y más granulares
    \item Tiempo de depuración significativamente reducido
    \item Mayor modularidad y menor acoplamiento
    \item Facilidad para mantenimiento y correcciones
\end{itemize}

\textbf{Resultado}: Arquitectura modular que facilitó el desarrollo posterior y la implementación de nuevas funcionalidades.

\section{Quinta iteración: Segundo intento de integración}
\label{sec:iteracion_5}

Con la nueva arquitectura modular, se realizó un segundo intento de integración que, aunque técnicamente superior, reveló la necesidad de cambios paradigmáticos más fundamentales.

\subsection{Análisis de la iteración}
\label{subsec:analisis_iteracion_5}

\textbf{Objetivos específicos}: Integrar módulos refactorizados mediante componente coordinador

\textbf{Problemas persistentes identificados}:
\begin{itemize}
    \item Comunicación entre módulos requería conocimiento específico interno
    \item Complejidad de configuración emergente
    \item Problemas de sincronización temporal no resueltos
\end{itemize}

\subsection{Diseño de la iteración}
\label{subsec:diseno_iteracion_5}

\textbf{Componente desarrollado}: Módulo integrador independiente que coordinaba ambos módulos

\textbf{Funcionalidades implementadas}:
\begin{itemize}
    \item Sistema de configuración unificado
    \item Coordinación visual de controles
    \item Exportación de datos separados
    \item \textit{Widgets} de ayuda para formatos de exportación
\end{itemize}

\subsection{Implementación y validación}
\label{subsec:implementacion_validacion_iteracion_5}

\textbf{Resultado}: Se logró integración visual y coordinación de inicio de grabaciones independientes desde el \acrshort{gui}, pero no se implementó acceso a \textit{streams} de datos para recolección y sincronización en tiempo real.

\textbf{Conclusión}: La integración funcionaba únicamente a nivel \acrshort{gui} sin capacidades de sincronización temporal real, motivando el cambio hacia arquitecturas \acrshort{api}.

\section{Sexta iteración: Arquitectura \acrshort{api} emergente}
\label{sec:iteracion_6}

Después de una pausa por otros compromisos académicos, el desarrollo se reanudó con una perspectiva nueva: implementar \acrshort{api} de acceso a datos internos para cada módulo.

\subsection{Análisis de la iteración}
\label{subsec:analisis_iteracion_6}

\textbf{Objetivos específicos}: Implementar acceso programático a datos internos mediante \acrshort{api} (RF-10).

\textbf{Requisitos abordados}: RF-10, RNF-04, RNF-06

\textbf{Casos de uso}: CU-03, CU-06

\subsection{Diseño de la iteración}
\label{subsec:diseno_iteracion_6}

\textbf{Arquitectura \acrshort{api} inicial}:
\begin{itemize}
    \item \acrshort{api} de acceso a datos (solo lectura)
    \item Sistema de suscripción a \textit{callbacks}
    \item Desarrollo independiente por módulo
\end{itemize}

\textbf{Evolución hacia explorador web}:
Debido a la necesidad de probar las \acrshort{api} desarrolladas, se implementó una aplicación web con \textit{FastAPI} que exponía los \textit{endpoints} internos. Se desarrollaron aplicaciones web separadas inicialmente, pero por mantenibilidad se refactorizó hacia un explorador web unificado.

\subsection{Implementación de la iteración}
\label{subsec:implementacion_iteracion_6}

\textbf{Desarrollo de \acrshort{api} de acceso}:
Se implementaron \acrshort{api} de solo lectura con \textit{callbacks} para el módulo de procesamiento de señales y sistema dirigido por eventos para el módulo de entrada.

\textbf{Explorador web unificado}:
Se desarrolló un módulo genérico donde los módulos independientes se registran dinámicamente, exponiendo sus métodos para visualización web. Incluye visualizadores tiempo real con \textit{Chart.js} y alternancia entre modos oscuro/claro.

\textbf{Visualizadores tiempo real}:
Para señales \acrshort{eeg}: gráficos con adaptación automática del eje Y, visualización separada/unificada, filtrado por tipo de canal. Para eventos de entrada: visualización \acrshort{json} estructurada específica del controlador.

\textbf{Problema crítico resuelto}: Se identificó que cerrar el visualizador \acrshort{gui} pausaba el \textit{stream} porque el único punto de \textit{polling} estaba en el visualizador. Se implementó \textit{polling} interno temporal en el colector.

\subsection{Validación de la iteración}
\label{subsec:validacion_iteracion_6}

\textbf{Pruebas realizadas}:
\begin{itemize}
    \item Verificación de \acrshort{api} de acceso a datos
    \item Validación de \textit{streaming} \textit{WebSocket}
    \item Pruebas de visualización web tiempo real
    \item Verificación de sistema de registro dinámico de módulos
\end{itemize}

\textbf{Resultado}: Arquitectura \acrshort{api} exitosa que estableció las bases para todo el desarrollo posterior.

\section{Séptima iteración: \acrshort{api} de control y extensiones}
\label{sec:iteracion_7}

Esta iteración completó la arquitectura \acrshort{api} con capacidades de control bidireccional y añadió soporte para controladores \acrshort{ps4}.

\subsection{Análisis de la iteración}
\label{subsec:analisis_iteracion_7}

\textbf{Objetivos específicos}: Implementar control remoto completo y soporte \acrshort{ps4} (RF-03, RF-10).

\textbf{Requisitos abordados}: RF-03, RF-10, RNF-04, RNF-07

\subsection{Diseño de la iteración}
\label{subsec:diseno_iteracion_7}

\textbf{\acrshort{api} de control}:
\begin{itemize}
    \item Sincronización bidireccional con \acrshort{gui} mediante señales \textit{Qt}
    \item Validación de parámetros y manejo de errores
    \item Operaciones \textit{thread-safe} con \textit{timeouts}
\end{itemize}

\textbf{Soporte \acrshort{ps4}}:
\begin{itemize}
    \item Integración \textit{PySDL2} para acceso de bajo nivel
    \item Inicialmente \textit{polling} con \textit{timers} \textit{Qt}
    \item Migración posterior a arquitectura dirigida por eventos
\end{itemize}

\subsection{Implementación de la iteración}
\label{subsec:implementacion_iteracion_7}

\textbf{Desarrollo del \acrshort{api} de control}:
Se implementó control remoto completo para el módulo de entrada con sincronización \acrshort{gui}-\acrshort{api} mediante señales \textit{Qt}. Se desarrolló plantilla inicial para el módulo \acrshort{eeg} basada en la estructura del módulo de entrada.

\textbf{Implementación del controlador \acrshort{ps4}}:
Se desarrolló con \textit{PySDL2} mediante conexión por cable. La implementación inicial usaba \textit{polling} con \textit{timers} \textit{Qt}, pero no funcionaba en modo \textit{headless} (\acrshort{cli} sin \textit{Qt}).

\textbf{Migración a arquitectura dirigida por eventos}:
Se cambió el colector \acrshort{ps4} a eventos usando hilo dedicado con detección de cambios por comparación de estados. Esto requirió optimizaciones por el alto volumen de eventos: \textit{batching}, \textit{workers}, \textit{buffer} circular, serialización con \textit{orjson}.

\textbf{Unificación de servidores}:
Se integraron \acrshort{api} de acceso y control en un único servidor web, registrando todos los módulos en lugar de servidores independientes.

\textbf{Tipado estático completo}:
Se implementó tipado con \textit{type hints} en todo el código y configuración de \textit{pre-commit hooks} para verificación automática.

\subsection{Validación de la iteración}
\label{subsec:validacion_iteracion_7}

\textbf{Pruebas realizadas}:
\begin{itemize}
    \item Verificación de control remoto completo
    \item Pruebas de sincronización \acrshort{gui}-\acrshort{api}
    \item Validación de controlador \acrshort{ps4} en modo eventos
    \item Verificación de optimizaciones de rendimiento
\end{itemize}

\textbf{Resultado}: Sistema de control bidireccional completo con soporte para dispositivos avanzados y optimizaciones de rendimiento.

\section{Octava iteración: Funcionalidades avanzadas \acrshort{eeg}}
\label{sec:iteracion_8}

La iteración final implementó capacidades avanzadas de procesamiento \acrshort{eeg} y completó la integración con bibliotecas científicas especializadas.

\subsection{Análisis de la iteración}
\label{subsec:analisis_iteracion_8}

\textbf{Objetivos específicos}: Integrar procesamiento avanzado y herramientas científicas (RF-06, RF-09).

\textbf{Requisitos abordados}: RF-06, RF-09, RNF-02, RNF-09

\textbf{Casos de uso}: CU-04, CU-05, CU-06

\subsection{Diseño de la iteración}
\label{subsec:diseno_iteracion_8}

\textbf{Componentes avanzados}:
\begin{itemize}
    \item Integración \textit{MNE-Python} para análisis neurocientífico
    \item Sistema \acrshort{ica} con procesamiento multiproceso
    \item Análisis \textit{wavelet} con múltiples familias
    \item Filtros digitales avanzados
\end{itemize}

\textbf{Arquitectura de servicios}:
\begin{itemize}
    \item Centralización en \textit{DataAccessService}
    \item Sistema de \textit{callbacks} optimizado
    \item Soporte para múltiples \textit{presets} de dispositivos
    \item \textit{Streaming} a archivo y red
\end{itemize}

\subsection{Implementación de la iteración}
\label{subsec:implementacion_iteracion_8}

\textbf{Finalización del \acrshort{api} de control \acrshort{eeg}}:
Se implementó control remoto completo con más de 40 señales \textit{Qt} para sincronización. Durante el desarrollo se identificaron numerosos problemas internos que se resolvieron, mejorando la robustez general.

\textbf{Integración \textit{MNE-Python}}:
Se desarrollaron servicios de conversión para transformar datos BrainFlow a formatos \textit{MNE} en tiempo real. Se identificó y corrigió el error crítico del \textit{get\_instance()} que causaba fallos generalizados.

\textbf{Centralización del acceso a datos}:
Se centralizó completamente el acceso en \textit{DataAccessService}, eliminando duplicaciones. Se migró de \textit{polling} a \textit{callbacks} (locales para baja latencia, \acrshort{api} directos o multiproceso).

\textbf{Algoritmos avanzados}:
Se implementó \acrshort{ica} con 1-16 componentes, procesamiento tiempo real y exportación de resultados. Las transformadas \textit{wavelet} incluyeron 15+ familias con configuración flexible. Los filtros avanzados expandieron las capacidades básicas (\textit{bandpass}, \textit{notch}, \textit{detrend}).

\textbf{Detección automática de canales}:
Se eliminó la restricción de asumir primeros 16 canales como \acrshort{eeg}. El sistema detecta automáticamente tipos de canales y crea pestañas especializadas con soporte para múltiples \textit{presets}.

\textbf{Streamers y corrección de grabación}:
Se implementó \textit{streaming} a archivo (datos procesados) y \textit{tcp}/\textit{udp} (datos \textit{RAW}, restringido a \textit{preset} 0). Se corrigieron problemas donde los filtros se aplicaban a canales de control, causando fallos en reproducción. Se migró la lógica de grabación a \textit{callbacks} eliminando pérdida de muestras.

\subsection{Validación de la iteración}
\label{subsec:validacion_iteracion_8}

\textbf{Pruebas realizadas}:
\begin{itemize}
    \item Validación de integración \textit{MNE-Python}
    \item Verificación de algoritmos \acrshort{ica} (funcionalidad, no corrección científica)
    \item Pruebas de transformadas \textit{wavelet} (sin errores, aplicación correcta)
    \item Validación de \textit{streaming} a archivo y red
    \item Verificación de grabación sin pérdida de muestras
\end{itemize}

\textbf{Métricas de rendimiento}:
\begin{itemize}
    \item Soporte validado hasta 1000Hz de frecuencia de muestreo
    \item Latencia de procesamiento < 100ms para algoritmos avanzados
    \item Pérdida de muestras < 1\% en grabación continua
    \item Tiempo de respuesta \acrshort{api} < 50ms promedio
\end{itemize}

\section{Implementación de interfaces de línea de comandos}
\label{sec:implementacion_cli}

Como complemento a las \acrshort{api} desarrolladas, se implementaron interfaces \acrshort{cli} completas para automatización y uso programático.

\subsection{Diseño de la arquitectura \acrshort{cli}}
\label{subsec:diseno_cli}

\textbf{Tecnologías utilizadas}: \textit{Typer}, \textit{Rich}

\textbf{Funcionalidades implementadas}:
\begin{itemize}
    \item Comandos de datos para acceso programático
    \item Comandos de interacción para control remoto
    \item Monitoreo en tiempo real con visualización \textit{Rich}
    \item Generación automática de ayuda con \textit{Typer}
\end{itemize}

\textbf{Problemas identificados}:
\begin{itemize}
    \item Logs de aplicación mezclados con salida \acrshort{cli}
    \item Dependencias \textit{Qt} ocultas impiden funcionamiento completamente \textit{headless}
    \item Problemas específicos en módulo de entrada (servicios sin inicializar)
    \item \acrshort{cli} de \acrshort{eeg} funciona correctamente excepto visualizadores avanzados
\end{itemize}

\section{Décima iteración: Intento de integración final}
\label{sec:decima_iteracion}

Se desarrolló un módulo de sincronización con \acrshort{lsl} que accedía a datos mediante \acrshort{api} internas con \textit{callbacks} locales, pero se pausó por problemas técnicos.

\subsection{Desarrollo del módulo de sincronización}
\label{subsec:modulo_sincronizacion}

\textbf{Diseño}: Módulo que accedía a datos de ambos módulos mediante \textit{callbacks} locales, integrando visualmente y coordinando operaciones.

\textbf{Problemas encontrados}:
\begin{itemize}
    \item Fallos en exportación de formatos \acrshort{edf} y \acrshort{csv}
    \item Problemas de compatibilidad y coherencia de datos
    \item Complejidad mayor de la anticipada
\end{itemize}

\textbf{Estado actual}: Desarrollo pausado debido a restricciones temporales. Los \textit{callbacks} locales funcionan correctamente y la integración conceptual está validada.

\section{Arquitectura final del sistema}
\label{sec:arquitectura_final}

El sistema desarrollado implementa una arquitectura modular basada en servicios con múltiples interfaces de acceso.

\subsection{Componentes principales}
\label{subsec:componentes_principales}

\textbf{Módulo de captura de entrada}:
\begin{itemize}
    \item Arquitectura modular con servicios especializados
    \item Soporte múltiple: teclado, ratón, controladores \acrshort{ps4}
    \item Sistema de eventos dirigido por eventos
    \item Interfaces: \acrshort{gui} \textit{PyQt5}, \acrshort{cli}, \acrshort{api} web
\end{itemize}

\textbf{Módulo de procesamiento \acrshort{eeg}}:
\begin{itemize}
    \item Arquitectura orientada a servicios
    \item Integración BrainFlow con múltiples fabricantes
    \item Procesamiento avanzado: \acrshort{ica}, \textit{wavelet}, \textit{MNE-Python}
    \item Visualización científica en tiempo real
\end{itemize}

\textbf{Explorador \acrshort{api}}:
\begin{itemize}
    \item \textit{Framework} web modular con \textit{FastAPI}
    \item Registro dinámico de módulos
    \item Visualización tiempo real con \textit{WebSocket}
    \item Documentación automática de \acrshort{api}
\end{itemize}

\subsection{Cumplimiento de requisitos}
\label{subsec:cumplimiento_requisitos}

\textbf{Requisitos funcionales}: Cumplimiento de RF-01 a RF-12, con RF-08 (sincronización en tiempo real) limitado a grabaciones independientes sincronizables \textit{offline}.

\textbf{Requisitos no funcionales}: Cumplimiento de RNF-01 a RNF-10 con limitación parcial en RNF-06 (dependencias \textit{Qt} en modo \textit{headless}).

\textbf{Casos de uso}: Soporte para CU-01 a CU-06, con CU-01 (generación de \textit{datasets}) requiriendo sincronización \textit{offline}.

\textbf{Capacidades del sistema}: Cada módulo puede grabar independientemente en \acrshort{csv} (también \acrshort{json} en el caso de entrada), permitiendo sincronización \textit{offline} posterior sin sobrecarga \acrshort{cpu} en tiempo real y evitando fallos de sincronización entre módulos.

El sistema desarrollado constituye una plataforma que integra capacidades tradicionalmente fragmentadas en el ecosistema \acrshort{bci}, proporcionando herramientas necesarias para investigación y desarrollo en interfaces cerebro-computadora.